\begin{obeactivity}
\textbf{Aktivitas 11.1: Simulasi Perhitungan Diffie-Hellman}

1. Gunakan parameter publik $p = 353$ dan $g = 3$.

2. Alice memilih kunci privat $a = 97$. Hitunglah nilai $A$.

3. Bob memilih kunci privat $b = 233$. Hitunglah nilai $B$.

4. Hitunglah kunci rahasia bersama $K$ yang dihasilkan oleh Alice dan Bob.

5. Diskusikan bagaimana jika $p$ yang digunakan adalah bilangan kecil. Apakah
   keamanan tetap terjamin?

6. Ulangi langkah 2 sampai 4 dengan perkakas simulasi daring yang memakai
   $G = 1601$ dan $N = 4789$. Bandingkan hasilnya dengan
   Contoh~\ref{ex:dh-1601}.
\end{obeactivity}

Langkah-langkah pada Aktivitas 11.1 dapat dijalankan dengan program pada
Listing~\ref{lst:diffie-hellman}. Program itu menghitung kunci publik dan kunci
bersama dari kedua pihak, sekaligus memeriksa apakah $g$ benar-benar akar
primitif dari $p$. Panggil fungsi \code{jalankan(p=353, g=3, a=97, b=233)}
untuk memperoleh hasil Contoh~\ref{ex:dh-353}.

\lstinputlisting[
    language=Python,
    firstline=10,
    lastline=58,
    firstnumber=1,
    caption={Pertukaran kunci Diffie-Hellman dan pemeriksaan akar primitif},
    label={lst:diffie-hellman}
]{\codefile{python/diffie_hellman.py}}

\begin{obeactivity}
\textbf{Aktivitas 11.2: Menelusuri Kriptografi Hibrida}

1. Tuliskan pasangan kunci milik Alice dan pasangan kunci milik Bob.

2. Ikuti lima langkah pada Subbagian~\ref{subsec:hibrida-eksplisit}. Tuliskan
   nilai yang dihitung Alice pada setiap langkah.

3. Sebutkan nilai yang melewati saluran dan nilai yang tidak pernah melewati
   saluran.

4. Jelaskan mengapa Bob dapat memperoleh $K$, sedangkan pihak lain tidak.
\end{obeactivity}

\begin{obeactivity}
\textbf{Aktivitas 11.3: Menelusuri Serangan Man-in-the-Middle}

1. Gambarkan posisi Alice, Mallory, dan Bob.

2. Tuliskan nilai yang dikirim Alice, nilai yang diganti Mallory, dan nilai yang
   diterima Bob.

3. Tuliskan $K_1$ dan $K_2$ yang dihitung oleh Alice, Bob, dan Mallory.

4. Jelaskan mengapa Mallory tidak perlu mengetahui $a$ atau $b$.
\end{obeactivity}

\begin{obeactivity}
\textbf{Aktivitas 11.4: Pertukaran Kunci Tiga Pihak}

1. Ikuti dua belas langkah pertukaran kunci tiga pihak pada Subbagian
   Pertukaran Kunci Tiga Pihak dengan lambang $a$, $b$, dan $c$.

2. Tunjukkan tiga jalur yang mencapai $g^{abc} \bmod p$.

3. Jelaskan mengapa ketiga jalur itu menghasilkan kunci yang sama.
\end{obeactivity}
